\paragraph{The two exchanges in the triangle}
Each block has length \(t\). The sum of the ranks at the two interval
nodes equals that of the two intermediate nodes, whose rank difference
is \(2t\). The numbers of negative signs satisfy the same additive
relation. The character for the exchange is the quotient of
\(\eta_i\xi^i\) and \(\eta_j\xi^j\). At the chosen places the
\((j-i)\)-th power of the fixed character is nontrivial, and at the twist
places all powers with \(1\le j-i\le m\) ramify.
Thus the points entering \(L_*\) in the exchange are \(1,\ldots,t\).
\leanlinks{\leansrc{New/Determinant/Chain/ChainNorms}{Checking the ranks, signs, and characters}}

To pass from point data to standard inclusions, work over the finite
residue ring \(R\). Let \(a',b'\) be the widths of the two completing
blocks and \(n'\) the height of the variable upper block. For a weight
\(p\) on \(n'\times a'\) matrices, use on the opposite side the datum
\[
 p_\mathrm{next}\begin{pmatrix}\eta&s\\0&I_{b'}\end{pmatrix}
     =p(-\eta)/|R|^{n'b'}
\]
and set it to zero elsewhere. With \(w=\binom{\eta}{0}\), the matrices
\((w\ \ I)\) and \(\binom{I}{-w}\) are the inverse-row and first-column
blocks of a common invertible completion. Counting the bases of the
two corresponding frames gives factors
\(|\mathrm{GL}_{n'+b'}(R)|\) and \(|\mathrm{GL}_{a'}(R)|\).
The upper shear \(\begin{psmallmatrix}I&s\\0&I\end{psmallmatrix}\) has
determinant \(1\) and preserves the kernel value. Average over \(s\) and
sum over \(\eta\) to obtain the stated weight.
For the reversed edge use the finite Fourier transform of the same
identity. Expanding it, the sums over the shear coordinates impose
the corresponding block-product conditions by additive orthogonality.
The remaining square block gives the matrix Gauss sum of a primitive
determinant character. Its modulus at the identity and the size of
the residue ring give the additional factors. Together with
\eqref{eq:exchanged-integrals}, this leaves \(L_*\) and modulus powers
as the twist-dependent absolute values for the adjoint exchange.
\leanlinks{\leansrc{New/Determinant/Chain/SquareLaw}{Shear averaging and finite Fourier transform for the inclusions}
  \leanlinksep \leansrc{New/Determinant/Chain/ChainNorms}{The two commutation relations in the triangle}}

\paragraph{The Gram form of a single lift}
The lift from rank \(r\) to rank \(r+t\) satisfies \(2r<r+t\).
The inner product formula in this range writes the product of the
two theta kernels as a single kernel on \(W_j\perp W_j^{\mathrm{opp}}\).
The Siegel--Weil expansion of this kernel is a sum over maximal
isotropic subspaces. The finite datum is supported where both blocks
of \(r\) columns are the specified primitive inclusions. Hence the
surviving isotropic subspaces project isomorphically onto both sides:
they are the graphs of rational isometries. Their zero slices are Weil
matrix coefficients. Unfolding the sum over graphs into the integral
over rational points gives the pairing of the forms with right
translations weighted by these coefficients.
\leanlinks{\leansrc{Lib/SiegelWeil/RallisInner/Doubling}{The doubled kernel and the graph terms}}

At the chosen level and the twist places the primitive inclusions
give the block datum. Its matrix coefficient integrates to the level
average times the mass of the datum, so on level-fixed forms the
average is the identity. The integral of the real matrix coefficient
is a fixed nonzero constant by reproduction for the corresponding
discrete series types. At the remaining finite places the unramified
vacuum coefficients give the translation series \(R_j\).
For a coset representative \(g\) supported away from the level and
twist places and their conjugates, its coefficient is
\((\chi_{V_j}(\det g)m_f(g)^{r+t})^{-1}\),
where \(m_f\) is the finite vacuum multiplier and \(\chi_{V_j}\) is the
character for the partner space.

The identity coset contributes the identity operator. Let \(a_w\)
be the count of neighbouring steps at an unramified place, weighted
by \(q_w^{-(r+t)/2}\). Since \(2r<r+t\), the places together with their
conjugates can be grouped so that the sum of weights in each group
is at most \(1/3\). All nonidentity terms in the series for such a
group have total coefficient modulus at most
\(\sum_{k\ge1}(1/3)^k=1/2\).
If the series sends a bounded \(f\) to zero, then
\(\|f\|_\infty\le\|f\|_\infty/2\), so \(f=0\).
Absolute convergence and factorization over disjoint sets of places
give injectivity of \(R_j\). If the lift of \(f\) is zero, use \(R_j f\)
as the test form in the Gram identity to obtain \(\|R_j f\|^2=0\).
This proves injectivity of each individual lift.
\leanlinks{\leansrc{Lib/SiegelWeil/RallisInner/SharpData}{Level averaging and the inner product formula}
  \leanlinksep \leansrc{Lib/SiegelWeil/UnramifiedSeries/RallisSeries}{Injectivity of the vacuum-coefficient series}}

The algebraic-form inner product is the Petersson product times the
level index and the dimension of the Weyl module. The index cancels
the denominator of the mass of the prescribed columns.
At a twist place \(w\) the projection onto the rank-\(r\) block and the
free block of width \(t\) give \((\mathrm Nw)^{rt}\).
Counting invertible blocks by
\(|\mathrm{GL}_n(\mathbb F_q)|=q^{n^2}\prod_{s=1}^n(1-q^{-s})\)
gives the remaining ratio \(\prod_{t<s\le t+r}(1-(\mathrm Nw)^{-s})\).
Thus the twist-dependent scalar is \(c_Q(\xi)\) below.
\leanlinks{\leansrc{New/Determinant/Norms/AdjointThetaLift}{The finite mass and the Weyl module dimension}}

To compare the good spaces on \(W_j\) and \(Y_0\), use their common
finite groups and infinite parameters. Only the signatures differ.
For each good label tuple there is a primitive label at a split
place above \(p_0\). Its central idempotent is attached to a cuspidal
character of the finite general linear group, so the local constant
terms along proper parabolics vanish. The trace comparison for
cuspidal test functions applies. At an auxiliary place, descend
to the stabilizer of a line and use its Hecke operator to impose a
simple zero of the characteristic polynomial over the residue field.
The determinant of this operator counts the spherical vectors and
returns the trace identity at the original level.
For a spherical Hecke operator \(T\), the traces of \(T^n\) coincide.
Newton's identities give equal characteristic polynomials and
determinants. Finitely many commuting semisimple operators separate
the eigensystems. Taking traces of interpolating polynomials that
select individual eigensystems gives equality of their multiplicities.
The sum \(R_j\) acts as one spherical Hecke operator on the two spaces,
so its determinant can be taken on the good space for \(Y_0\).
\leanlinks{\leansrc{Lib/Adelic/TraceFormula/CuspidalComparison}{The line stabilizer and comparison for cuspidal test functions}
  \leanlinksep \leansrc{New/Spanning/Multiplicities/EqualMultiplicities}{Equality of Hecke eigensystem multiplicities}}
